% =============================================================================
% A2 MANUSCRIPT DRAFT — assembled 2026-07-06 from the accepted final data lock
% [local path omitted]
% Claim wording is bound VERBATIM to locked/A2_ALLOWED_CLAIMS.md.
% Locked inputs enter via \input from locked/ (byte-identical, hash-verified).
% Final draft: theory-pass citation/context slots filled; locked results/tables remain unchanged.
% =============================================================================
\documentclass[11pt]{article}
\usepackage{amsmath,amssymb}
\usepackage{graphicx}
\usepackage[margin=1in]{geometry}
\usepackage{booktabs}
\usepackage{url}
\usepackage{pdflscape}

\title{Surface-independent stability of equilateral and squeezed bispectra\\
on the frozen BIH-VI benchmark: a bounded-$\eta_H$ all-operator production scan}
\author{Lou Thomas\\ \small Independent researcher; Boundary-Inversion Cosmology project}
\date{July 6, 2026 --- draft assembled from locked production data}

\begin{document}
\maketitle

\begin{abstract}
We report a production-scale stability scan of the primordial curvature bispectrum on the frozen
BIH-VI benchmark background, an effective-field-theory (EFT) clock realization of an
entropy-selected density-shield scenario. The scan applies a validated bounded-$\eta_H$
all-operator instrument --- a finite-basis cubic-action ledger with sealed kernels, independent
unit-test engines, and declared hard gates --- to equilateral, squeezed
($q = 10^{-3}$), and folded triangle configurations over a hard grid $K \in [0.010, 0.060]$ and a
provisional grid $K \in \{0.072, 0.090, 0.119\}$. All 13 declared production gates pass. In the
claim-bearing hard envelope $K \leq 0.06$, equilateral and squeezed configurations are stable under
the sealed surface-independence (G-S) criteria, with $f_{\rm NL}$ in
$[0.02230, 0.03822]$ (equilateral) and $[0.05507, 0.05906]$ (squeezed) at the reference surface.
Folded/flattened configurations are diagnostic/report-only under the current sealed instrument; in
this run, no folded/flattened row passed the G-FLAT diagnostic. Every reported number is bound to a
cryptographically hashed artifact chain with two-party verification; the complete locked data
package accompanies this paper. This is an EFT-level methodology and stability result; no
UV-completion claim is made.
\end{abstract}

\section{Introduction}
Stability of a computed observable under changes of bookkeeping choices --- evaluation surface,
regulator resolution, operator-basis completion --- is a precondition for treating that observable
as physical rather than artifactual. This paper reports such a stability program, executed to
completion, for the tree-level curvature bispectrum on a single frozen benchmark background of the
BIH-VI construction. The calculation is deliberately framed within the standard language of
single-clock inflationary perturbation theory and the effective field theory of inflation
\cite{maldacena2003,cheung2008}, where the curvature perturbation is the observable clock fluctuation
and cubic interactions must be combined with the appropriate constraint and field-redefinition terms.
Sharp time-dependent features are known to create localized oscillatory signatures in correlation
functions and to stress numerical in-in implementations \cite{chen2007,adshead2012}. The present work
uses that setting as a methodology test: a sharp exit transition is allowed, but every reported
bispectrum number must survive operator-completion, normalization, convergence, and surface-independence
gates before it is treated as claim-bearing.

The background motivation for BIH-VI is an entropy-selected apparent-horizon density shield, in the
spirit of horizon-thermodynamic approaches to gravitational dynamics \cite{jacobson1995,hayward1998}.
For observational orientation, the frozen benchmark was chosen from a Planck-compatible scalar/tensor
parameter island \cite{planck2018params}. These motivations do not enter as additional claims in this
paper: the scan below tests only the stated EFT-clock benchmark and its bounded-$\eta_H$ bispectrum
instrument.

The working discipline throughout is that frame- or scheme-dependent descriptions must be reduced
to invariant, operationally measurable statements before they count as physics. Concretely: every
claim in this paper is (i) evaluated at multiple hypersurfaces and gated on surface independence,
(ii) computed from a sealed operator ledger whose kernels are verified against an independent
engine, and (iii) bound to hash-locked artifacts so that each number is reproducible from the
archived instrument.

Three boundaries are declared at the outset. First, the BIH-VI benchmark is an EFT-level
realization and phenomenological test set, not a final ultraviolet-complete microscopic theory; no
UV-completion claim is made. Second, the claim-bearing results are confined to the hard envelope
$K \leq 0.06$ in the equilateral and squeezed families; no all-shape closure is claimed.
$K \leq 0.06$ is the current claim-bearing boundary of the sealed A2 instrument for the frozen
benchmark, not a physical cliff in the EFT-clock model. Third,
folded/flattened configurations are reported under a diagnostic-only disposition (Decision~A,
Section~\ref{sec:production}); no physical conclusion --- neither stability nor instability --- is
drawn from them.

\subsection*{Nomenclature and definitions}
\begin{description}
\item[hard envelope] The claim-bearing set: equilateral and squeezed hard-grid claim rows with
  $K \leq 0.06$.
\item[$K$ (grid label)] The production-grid label used throughout this paper (the per-mode scale;
  for equilateral rows $k_1 = k_2 = k_3 = K$). The triangle perimeter
  $K_\Sigma = k_1 + k_2 + k_3$ is not used as a variable in this paper; where a perimeter is
  needed it is written explicitly.
\item[sealed G-S criteria] The pre-registered G-S surface-independence, frozen-window, Wronskian,
  convergence, and continuity-sentinel checks, fixed before the production run.
\item[claim-bearing rows] Rows included in \texttt{claim\_rows.csv} and authorized by
  \texttt{A2\_ALLOWED\_CLAIMS.md}.
\item[non-claim-bearing rows] Rows reported for diagnostics, provisional context, or flagged
  status but excluded from the A2 claim.
\item[folded/flattened rows] Diagnostic/report-only unless G-FLAT explicitly passes; in this run,
  no folded/flattened row passed G-FLAT.
\end{description}

\section{Benchmark background}
\label{sec:benchmark}
The frozen BIH-VI benchmark is specified by
\begin{equation}
H^2 = H_c^2\, y(2-y), \qquad
y_N = -m(y)\, y, \qquad
m(y) = \mu (1-y)^p + m_{\rm post}\, S(y),
\end{equation}
with switch profile
\begin{equation}
S(y) = \frac{1}{1 + \exp\!\big[(y - y_g)/\Delta\big]},
\end{equation}
and parameters
\begin{equation}
p = 3,\quad \mu = 1,\quad c_s = 1,\quad N_* = 55,\quad y_g = \tfrac{2}{3},\quad
\Delta = 0.01,\quad m_{\rm post} = 4,\quad H_c = 3.8\times 10^{-6}\, M_{\rm Pl}.
\end{equation}
Numerics use the dimensionless convention $H(y_g) = 1$; the value of $H_c/M_{\rm Pl}$ enters only
as provenance for physical-amplitude conversions. ($c_s = 1$ is implicit in the mode equation of
the sealed instrument; $c_s$ and $N_*$ are fixed by the frozen benchmark specification,
hash-referenced in Appendix~\ref{app:repro}.) The background is integrated at
$\Delta N = 10^{-3}$ over $N \in [-22, 0.45]$ with fully analytic $\eta$ and $\eta_N$ (no numerical
differentiation through the switch), and the slow-roll functions $(\epsilon, \eta)$ follow from
$m(y)$ as in the sealed instrument.

\section{Instrument: the bounded-$\eta_H$ all-operator ledger}
\label{sec:instrument}
The bispectrum is assembled from the sealed finite-basis ledger
\begin{equation}
B_{\rm total} = O_1 + O_2 + O_{3a} + O_{3b} + O_{3c} + V_1 + V_2 + f_2 + f_3 + f_4,
\end{equation}
where $O_1 = a^3 \epsilon^2 \zeta \dot\zeta^2$,
$O_2 = a\,\epsilon^2 \zeta (\partial\zeta)^2$,
$O_{3a} = -2 a\,\epsilon\, \dot\zeta\, \partial\zeta \cdot \partial\psi$,
$O_{3b} = (\epsilon/2a)\, \partial\zeta \cdot \partial\psi\, \partial^2\psi$,
$O_{3c} = (\epsilon/4a)\, \partial^2\zeta\, (\partial\psi)^2$,
$V_{1,2}$ are the bounded-$\eta$ replacement kinetic and gradient terms, and $f_{2,3,4}$ are the
nonlocal/decaying field-redefinition contributions. The amplitude is normalized as
\begin{equation}
f_{\rm NL} = \frac{5}{6}\, \frac{B_{\rm total}}{P_1 P_2 + P_1 P_3 + P_2 P_3},
\end{equation}
with dimensionful power spectra $P_i = |\zeta_{k_i}|^2$; dimensionless $\Delta^2$ combinations are
excluded from all denominators by construction. The following ledger lines are forbidden and absent:
$O_5$, standalone boundary terms ($B_{\rm bdry}$) \cite{arroja2011}, standalone equation-of-motion
terms ($B_{\rm EOM}$), the local $\zeta^2$ field-redefinition line, and any kernel missing its
assignment sum. Raw $O_4$ never enters any physical row or production observable: it is entirely absent from
the sealed production script, and in the certified lineage instrument it appears only as an
internal diagnostic (the synthetic-context kernel unit test and the regulated IBP-identity gate),
as permitted.

Mode functions are integrated with an RK4 scheme from deep inside the horizon
($k/aH \geq 60$) with Bunch--Davies initialization, and bulk integrals carry the sealed
$i\epsilon$ regulator ($\epsilon_{\rm reg} = 0.03$). Kernel implementations are verified against an
independent vector/bijection engine (gates UT-K, UT-FR, tolerance $10^{-10}$), and the Wronskian
identity is enforced along every mode (gate G-W).

\section{Validation lineage}
\label{sec:lineage}
The instrument used here is sealed: it is the same code path certified by two prior gate programs,
executed and archived under a hash-locked reproduction protocol (Appendix~\ref{app:repro}).
\begin{itemize}
\item \textbf{G-S surface-independence sweep (v3).} The full ledger was evaluated at five
      hypersurfaces $N \in \{0.20, 0.25, 0.30, 0.35, 0.40\}$ for squeezed and equilateral rows.
      Per-row spread gates, a vacuous-pass guard, smallest-$K$ asymptotic requirements, and a
      drift-scaling slope test were all declared in advance; the certified run passed with six
      primary rows and established the empirical frozen envelope $K \leq 0.06$ used here.
\item \textbf{BENCH02 flat-gauge cross-check.} An independent ADHL flat-gauge implementation passed
      11 of 11 declared gates against the same benchmark, closing the residual $v_3$ difference and
      certifying the template, background, phase, envelope, squeezed, and convergence behavior of
      the instrument.
\item \textbf{Production continuity.} The production run verifies continuity sentinels against the
      certified G-S reference values at $K \in \{0.020, 0.040, 0.060\}$ for both claim families,
      plus a $\Delta N/2$ convergence sentinel, inside the same gate set
      (Table~\ref{tab:gates}).
\end{itemize}

\section{Production scan design}
\label{sec:production}
The production configuration scans three triangle families --- equilateral
($k_1 = k_2 = k_3 = K$), squeezed ($q = k_L/k_S = 10^{-3}$), and folded
(near-degenerate, $\delta = 10^{-2}$ with a $\delta = 10^{-3}$ tightening probe) --- over the hard
grid $K \in \{0.010, 0.015, \ldots, 0.060\}$ and the provisional grid
$K \in \{0.072, 0.090, 0.119\}$. Thirteen gates are declared: the instrument gates (UT-K, UT-FR,
BG-ETA, G-W, CONV), the G-S continuity sentinels, per-family hard-row stability gates, the folded
disposition gates (folded rows either pass G-FLAT explicitly or are flagged; no folded claim
without a G-FLAT pass), claim-hygiene gates (no provisional claim; no $K > 0.06$ claim), and the
execution lock. Per-family row accounting is summarized in Table~\ref{tab:families}.

\textbf{Decision A (folded/flattened disposition).} Folded/flattened configurations are
diagnostic/report-only under the current sealed instrument unless G-FLAT explicitly passes. The
G-FLAT diagnostic probes each folded row three ways: regulator halving
($\epsilon_{\rm reg} \to \epsilon_{\rm reg}/2$), flattening-parameter tightening
($\delta = 10^{-2} \to 10^{-3}$), sign-flip and cancellation-factor screens, in addition to the
standard surface-stability and freeze requirements. This disposition was adopted after an earlier
wide-grid exploratory scan showed that, near the flattened boundary at the edge of an extended
$K$ grid, the ledger sum passes through a zero: the total becomes a small residual of large
cancelling operator contributions, and log-slope estimators straddling the sign change report
spuriously large scaling exponents. That behavior is an instrument-scope limitation of the sealed
kernels in the flattened limit --- diagnosed, bounded, and archived --- not a physical instability
claim; resolving it is deferred to a separately validated instrument revision.

The production scan executed exactly once, under an execution lock requiring explicit
authorization flags, from a hash-verified script, with the frozen-artifact record confirmed before
and after execution.

% ===================== LOCKED RESULTS SECTION (verbatim) =====================
\section{Results}

\subsection{Production Run State}

The validated bounded-$\eta_H$ all-operator A2 instrument (script SHA-256 \texttt{35fd365916395f17\ldots}) was applied once to the frozen BIH-VI benchmark under Decision~A. The production run (\texttt{A2\_PRODUCTION\_LOCAL}) completed with verdict PASS in 7.280172109603882~s, scanning 42 shape rows over the hard grid $K \in [0.01, 0.06]$ (step $0.005$) and the provisional grid $K \in \{0.072, 0.090, 0.119\}$ for the equilateral, squeezed ($q = 10^{-3}$), and folded families.

\subsection{Gate Summary}

All 13 declared gates passed (Table~\ref{tab:gates}): kernel unit tests (UT-K, UT-FR $\leq 10^{-10}$), background consistency (BG-ETA), the Wronskian identity (G-W), $\Delta N/2$ convergence (CONV), continuity sentinels against the certified G-S reference values, per-family hard-row stability, and the four claim-hygiene gates (no folded claim without G-FLAT, no provisional claim, no $K > 0.06$ claim, execution lock).

\subsection{Claim-Bearing Rows}

The claim-bearing set (Table~\ref{tab:claims}) comprises 22 rows: 11 equilateral and 11 squeezed hard-envelope configurations with $K \leq 0.06$, each frozen at all five evaluation surfaces and stable under the sealed G-S criteria. Over this envelope $f_{\rm NL}$ spans $[0.02230, 0.03822]$ (equilateral, mean $0.03146$) and $[0.05507, 0.05906]$ (squeezed, mean $0.05769$).

\subsection{Folded/Flattened Diagnostics}

Folded/flattened configurations are diagnostic/report-only under the current sealed instrument (Decision~A). Of 14 G-FLAT diagnostic evaluations, 0 passed; all folded rows are therefore reported with full diagnostics (Table~\ref{tab:gflat}) and excluded from claims. No physical conclusion is drawn from flagged folded rows.

\subsection{Scope and Forbidden Claims}

The validated bounded-$\eta_H$ all-operator A2 instrument was applied to the frozen BIH-VI benchmark. In the claim-bearing hard envelope K $\leq$ 0.06, equilateral and squeezed configurations are stable under the sealed G-S criteria. Folded/flattened configurations are diagnostic/report-only under the current sealed instrument; in this run, no folded/flattened row passed G-FLAT.

This is an EFT-level A2 production scan; no UV-completion claim is made. The following claims are explicitly not made: full all-shape closure; folded/flattened physical instability; BIH-VI observational detection; Hubble-tension solution.

% =============================================================================

\section{Scope, limitations, and outlook}
\label{sec:scope}
The allowed claim of this paper, in full, is: \emph{The validated bounded-$\eta_H$ all-operator A2
instrument was applied to the frozen BIH-VI benchmark. In the claim-bearing hard envelope
$K \leq 0.06$, equilateral and squeezed configurations are stable under the sealed G-S criteria.
Folded/flattened configurations are diagnostic/report-only under the current sealed instrument; in
this run, no folded/flattened row passed G-FLAT.}

The following claims are explicitly \emph{not} made: full all-shape closure; folded/flattened
physical instability; BIH-VI observational detection; UV completion; Hubble-tension solution.
Nothing in this paper uses parked interpretive material as evidence. The results are confined to
one frozen benchmark background and to the stated envelope; extension of the claim envelope, or a
claim-bearing treatment of the flattened limit, requires a separately validated instrument revision
and its own gate program.
No interpretive C4/C5/C6 material is used in this manuscript. Such material remains outside the
load-bearing A2 analysis and is reserved for separate future work.

This audit is specific to the frozen BIH-VI EFT-clock benchmark, the bounded-$\eta_H$ all-operator
ledger, and the declared A2 regulator/validation contract. Comparisons with other sharp-feature
calculations require matching conventions, basis choices, field redefinitions, power-spectrum
normalization, and regulator prescriptions before numerical discrepancies are interpreted
physically.

\appendix
\section{Reproducibility and artifact chain}
\label{app:repro}
Every number in this paper derives from a hash-locked artifact chain, archived in an append-only
ledgered hub and verified by two independent parties (generation-side audit and an external
reviewer recomputation, agreeing on all hashes). The chain is:
\begin{center}
\begin{tabular}{lll}
\toprule
artifact & role & SHA-256 (truncated) \\
\midrule
benchmark handoff specification & frozen benchmark definition & \texttt{ceba4ebc\ldots7288fad2} \\
\texttt{a2\_gs\_surface\_sweep\_v3.py} & certified G-S instrument & \texttt{3289ccb4\ldots4e48daed} \\
\texttt{bench02\_adhl\_flat.py} & certified flat-gauge cross-check & \texttt{b1071166\ldots ac4fa4646} \\
\texttt{a2\_production\_scan.py} & sealed production script & \texttt{35fd3659\ldots58c78aae8} \\
production output package (zip) & primary run record & \texttt{e8500cc3\ldots d2251c266} \\
final data lock package (zip) & frozen, reviewed data release & \texttt{061dc918\ldots8372cfbe1a} \\
\bottomrule
\end{tabular}
\end{center}
Full values:
\begin{flushleft}\footnotesize\ttfamily
spec:\ ceba4ebce5a5ce9d7fc7eaf44ff35f304f1c91cb1e079eb67df532347288fad2\\
gs\_v3:\ 3289ccb4d7c0443e335c4b7d01d40946770a484806eb713cf33bed394e48daed\\
bench02:\ b10711663374cd7e269bec7f83b3b1847e353eb925bd67f46d4f4f5ac4fa4646\\
production\_scan:\ 35fd365916395f171fee9278480bd102b267f545fb389d6039cfe0058c78aae8\\
production\_zip:\ e8500cc3db9152200a9ab6242bb22c35c26589976425ccfa978e171d2251c266\\
data\_lock\_zip:\ 061dc91860f355355fc503d36571e4bf571d70dcacdf1aa0c531fa8372cfbe1a
\end{flushleft}
Full 64-hexit hashes for every file \emph{inside} the data release, including all tables printed
here, are recorded in \texttt{SHA256SUMS.txt} and \texttt{A2\_ARTIFACT\_MANIFEST.csv} within the
final data lock package, which constitutes the data-availability release for this paper. The
lineage-artifact and package-level hashes above are recorded in the hub's append-only audit
ledgers and in the provenance record accompanying this manuscript; the production run wrote its
own ledger line at completion.

\section{Tables}
\begin{table}[htbp]
\centering
\caption{Production gate summary (13 gates).}
\label{tab:gates}
\begin{tabular}{lll}
\hline
gate & status & key value \\
\hline
UT\_K & PASS & value=2.66e-15 \\
UT\_FR & PASS & value=1.09e-15 \\
BG\_ETA & PASS & etaN\_grad\_rel=8.7e-05 \\
G\_W & PASS & worst\_surface\_region=2.15e-06 \\
CONV & PASS & rel\_dev=0.000841 \\
G\_S\_CONTINUITY\_SENTINELS & PASS &  \\
EQUILATERAL\_HARD\_ROWS\_STABLE & PASS & n=11 \\
SQUEEZED\_HARD\_ROWS\_STABLE & PASS & n=11 \\
FOLDED\_GFLAT\_OR\_FLAGGED & PASS & n=14 \\
NO\_FOLDED\_CLAIM\_WITHOUT\_GFLAT & PASS &  \\
NO\_PROVISIONAL\_CLAIM & PASS &  \\
NO\_K\_GT\_0P06\_CLAIM & PASS &  \\
EXECUTION\_LOCK & PASS &  \\
\hline
\end{tabular}
\end{table}

\begin{table}[htbp]
\centering
\caption{Claim-bearing rows: equilateral and squeezed hard envelope, $K \leq 0.06$.}
\label{tab:claims}
\begin{tabular}{llrrrrl}
\hline
family & shape\_id & $K$ & $f_{\rm NL}$@0.35 & $f_{\rm NL}$ spread & $B$ drift & status \\
\hline
equilateral & equilateral\_K0.010 & 0.010 & 0.0382177 & 2.272e-05 & 5.744e-04 & claim \\
equilateral & equilateral\_K0.015 & 0.015 & 0.0361488 & 5.104e-05 & 1.367e-03 & claim \\
equilateral & equilateral\_K0.020 & 0.020 & 0.0370021 & 9.068e-05 & 2.369e-03 & claim \\
equilateral & equilateral\_K0.025 & 0.025 & 0.034999 & 1.415e-04 & 3.914e-03 & claim \\
equilateral & equilateral\_K0.030 & 0.030 & 0.0344472 & 2.035e-04 & 5.720e-03 & claim \\
equilateral & equilateral\_K0.035 & 0.035 & 0.0321143 & 2.767e-04 & 8.350e-03 & claim \\
equilateral & equilateral\_K0.040 & 0.040 & 0.0309027 & 3.609e-04 & 1.132e-02 & claim \\
equilateral & equilateral\_K0.045 & 0.045 & 0.0290793 & 4.562e-04 & 1.522e-02 & claim \\
equilateral & equilateral\_K0.050 & 0.050 & 0.0259711 & 5.624e-04 & 2.102e-02 & claim \\
equilateral & equilateral\_K0.055 & 0.055 & 0.0249104 & 6.798e-04 & 2.647e-02 & claim \\
equilateral & equilateral\_K0.060 & 0.060 & 0.0223037 & 8.079e-04 & 3.512e-02 & claim \\
squeezed & squeezed\_q0.001\_K0.010 & 0.010 & 0.0573306 & 1.857e-05 & 3.016e-04 & claim \\
squeezed & squeezed\_q0.001\_K0.015 & 0.015 & 0.0583783 & 4.173e-05 & 6.644e-04 & claim \\
squeezed & squeezed\_q0.001\_K0.020 & 0.020 & 0.0588645 & 7.408e-05 & 1.169e-03 & claim \\
squeezed & squeezed\_q0.001\_K0.025 & 0.025 & 0.0590595 & 1.156e-04 & 1.816e-03 & claim \\
squeezed & squeezed\_q0.001\_K0.030 & 0.030 & 0.0588911 & 1.662e-04 & 2.619e-03 & claim \\
squeezed & squeezed\_q0.001\_K0.035 & 0.035 & 0.0585432 & 2.259e-04 & 3.581e-03 & claim \\
squeezed & squeezed\_q0.001\_K0.040 & 0.040 & 0.0581061 & 2.946e-04 & 4.706e-03 & claim \\
squeezed & squeezed\_q0.001\_K0.045 & 0.045 & 0.057511 & 3.723e-04 & 6.011e-03 & claim \\
squeezed & squeezed\_q0.001\_K0.050 & 0.050 & 0.0568472 & 4.590e-04 & 7.499e-03 & claim \\
squeezed & squeezed\_q0.001\_K0.055 & 0.055 & 0.0559888 & 5.546e-04 & 9.204e-03 & claim \\
squeezed & squeezed\_q0.001\_K0.060 & 0.060 & 0.0550745 & 6.591e-04 & 1.113e-02 & claim \\
\hline
\end{tabular}
\end{table}

% Wide locked tables are placed on landscape pages; locked file contents are untouched.
\clearpage
\begin{landscape}
\footnotesize
\begin{table}[htbp]
\centering
\caption{Per-family production summary.}
\label{tab:families}
\begin{tabular}{lrrrrrrrr}
\hline
family & rows\_tested & hard\_rows & provisional\_rows & claim\_rows & flagged\_rows & worst\_fNL\_spread & worst\_B\_spread\_rel & max\_abs\_fNL\_ref035 \\
\hline
equilateral & 14 & 11 & 3 & 11 & 3 & 0.003127 & 0.2604 & 0.03822 \\
squeezed & 14 & 11 & 3 & 11 & 3 & 0.00255 & 0.05894 & 0.05906 \\
folded & 14 & 11 & 3 & 0 & 14 & 0.002033 & 0.02211 & 0.09448 \\
\hline
\end{tabular}
\end{table}

\clearpage
\tiny
\setlength{\tabcolsep}{2pt}
\begin{table}[htbp]
\centering
\caption{Folded/flattened G-FLAT diagnostics (report-only under Decision A).}
\label{tab:gflat}
\begin{tabular}{lrrrrrrrlll}
\hline
shape\_id & K & fNL\_ref035 & fNL\_ieps\_half & fNL\_delta\_tight & instrument\_dev & delta\_dev & cancellation\_factor & sign\_flip & gflat\_pass & gflat\_reason \\
\hline
folded\_delta0.01\_K0.010 & 0.01 & 0.08988 & 0.275 & 0.0759 & 0.6732 & 0.1555 & 1.602 & False & False & flat\_limit \\
folded\_delta0.01\_K0.015 & 0.015 & 0.09192 & 0.2818 & 0.07751 & 0.6738 & 0.1567 & 1.669 & False & False & flat\_limit \\
folded\_delta0.01\_K0.020 & 0.02 & 0.09144 & 0.2756 & 0.0788 & 0.6682 & 0.1382 & 1.777 & False & False & flat\_limit \\
folded\_delta0.01\_K0.025 & 0.025 & 0.09324 & 0.2849 & 0.07723 & 0.6727 & 0.1718 & 1.829 & False & False & flat\_limit \\
folded\_delta0.01\_K0.030 & 0.03 & 0.0929 & 0.2822 & 0.07993 & 0.6708 & 0.1396 & 1.926 & False & False & flat\_limit \\
folded\_delta0.01\_K0.035 & 0.035 & 0.09404 & 0.2894 & 0.08158 & 0.6751 & 0.1325 & 1.988 & False & False & flat\_limit \\
folded\_delta0.01\_K0.040 & 0.04 & 0.09358 & 0.2876 & 0.08049 & 0.6746 & 0.1399 & 2.086 & False & False & flat\_limit \\
folded\_delta0.01\_K0.045 & 0.045 & 0.09347 & 0.2884 & 0.08017 & 0.6759 & 0.1423 & 2.177 & False & False & flat\_limit \\
folded\_delta0.01\_K0.050 & 0.05 & 0.09448 & 0.2965 & 0.07786 & 0.6814 & 0.1759 & 2.238 & False & False & flat\_limit \\
folded\_delta0.01\_K0.055 & 0.055 & 0.09301 & 0.29 & 0.07943 & 0.6792 & 0.146 & 2.367 & False & False & flat\_limit \\
folded\_delta0.01\_K0.060 & 0.06 & 0.09301 & 0.2927 & 0.07958 & 0.6823 & 0.1443 & 2.457 & False & False & flat\_limit \\
folded\_delta0.01\_K0.072 & 0.072 & 0.09321 & 0.302 & 0.07646 & 0.6914 & 0.1797 & 2.672 & False & False & flat\_limit \\
folded\_delta0.01\_K0.090 & 0.09 & 0.08963 & 0.2951 & 0.07585 & 0.6962 & 0.1537 & 3.144 & False & False & flat\_limit \\
folded\_delta0.01\_K0.119 & 0.119 & 0.0837 & 0.2882 & 0.06904 & 0.7096 & 0.1752 & 4.032 & False & False & flat\_limit \\
\hline
\end{tabular}
\end{table}

\setlength{\tabcolsep}{6pt}
\end{landscape}
\clearpage


\begin{thebibliography}{99}
\bibitem{maldacena2003}
J.~M. Maldacena,
``Non-Gaussian features of primordial fluctuations in single field inflationary models,''
\emph{Journal of High Energy Physics} \textbf{2003}(05), 013 (2003).

\bibitem{cheung2008}
C.~Cheung, P.~Creminelli, A.~L. Fitzpatrick, J.~Kaplan, and L.~Senatore,
``The effective field theory of inflation,''
\emph{Journal of High Energy Physics} \textbf{2008}(03), 014 (2008).

\bibitem{chen2007}
X.~Chen, R.~Easther, and E.~A. Lim,
``Large non-Gaussianities in single field inflation,''
\emph{Journal of Cosmology and Astroparticle Physics} \textbf{2007}(06), 023 (2007).

\bibitem{adshead2012}
P.~Adshead, C.~Dvorkin, W.~Hu, and E.~A. Lim,
``Non-Gaussianity from step features in the inflationary potential,''
\emph{Physical Review D} \textbf{85}, 023531 (2012).

\bibitem{arroja2011}
F.~Arroja and T.~Tanaka,
``A note on the role of the boundary terms for the non-Gaussianity in general k-inflation,''
\emph{Journal of Cosmology and Astroparticle Physics} \textbf{2011}(05), 005 (2011).

\bibitem{jacobson1995}
T.~Jacobson,
``Thermodynamics of spacetime: The Einstein equation of state,''
\emph{Physical Review Letters} \textbf{75}, 1260--1263 (1995).

\bibitem{hayward1998}
S.~A. Hayward,
``Unified first law of black-hole dynamics and relativistic thermodynamics,''
\emph{Classical and Quantum Gravity} \textbf{15}, 3147--3162 (1998).

\bibitem{planck2018params}
Planck Collaboration,
``Planck 2018 results. VI. Cosmological parameters,''
\emph{Astronomy \& Astrophysics} \textbf{641}, A6 (2020).

\end{thebibliography}


\end{document}
